\begin{statement}{P3K.48}{Proved}
The actual selected-cusp family of P3F.54/P3F.56 admits a periodic
orientation-preserving ordinary smooth product of its labelled
three-curve plumbing, based at its initial fibre and retaining its
two labelled vertices and complex curve orientations.
In the explicitly constructed product, after the actual based
unit-logarithm contraction, the physical circle on the fixed-curve
normal coordinate has axis
\[
 \nu(s,\theta)=H(s,\theta)iH(s,\theta)^{-1},\qquad
 \deg\nu=+3,
\]
where the domain is the cylinder with each end circle collapsed,
oriented by \(ds\wedge d\theta\), and the target sphere has its
outward orientation. The closed fixed-curve normal-frame comparison has
factor degrees \(6,0\) in complex-curve-times-parameter orientation,
and \((-6,0)\) in the Hopf-sphere orientation of P3K.23.
The full-unit circle action is displayed in
the proof. This product is not delta-equivariant. No circle-trace
correction, common-exterior variation matrix, nonfixed relative
integer or additional fan return is assigned here.
\end{statement}
\begin{proof}
\par

\medskip\noindent\textbf{OP.1. One actual vertex product and all three normal transition loops}

\par\nobreak

Use PO.1 of P3K.47 and its ordered charts \((a,b,c)\), \((u,v,w)\).
Keep the first vertex coordinates and use at the second vertex the
actual unit map
\[
                  (u,v,w)\longmapsto(u,\eta v,\kappa w).
                                                               \tag{OP.1}
\]
Here the units are first restricted to the curves, hence to \(q=0\).
The higher terms of their original chart maps are restored below.
The reference vertex at \(\rho_*\) uses the same map with its initial
constant units. This makes (OP.1) an actual based vertex comparison,
rather than a replacement of the initial coordinate system.

The tangent coordinates on each pair of curve charts now have literal
reciprocal transition. With the following ordered complex normal
coordinates the three transitions are
\[
 \begin{array}{c|c|c|c}
  &z&\text{first normal pair}&\text{second normal pair}\\ \hline
 C_0&b&(a,c)&(\eta v,u)\\
 C_1&c&(a,b)&(\kappa w,u)\\
 C_2&a&(b,c)&(\eta v,\kappa w).
 \end{array}
\]
Their normal transition matrices are \(zD_i(\rho)\), where
\[
 D_0=\operatorname{diag}(\eta/\alpha,\kappa\alpha),\quad
 D_1=\operatorname{diag}(\kappa/g,g\eta),\quad
 D_2=\operatorname{diag}(\eta,\kappa).                 \tag{OP.2}
\]
These are direct substitutions in the three rows of PO.1. Put
\(K_i=D_i(\rho_*)^{-1}D_i(\rho)\). In particular
\[
 \begin{aligned}
 K_0&=\operatorname{diag}
   ((\eta/\eta_*)(\alpha_*/\alpha),
             (\kappa/\kappa_*)(\alpha/\alpha_*)),\\
 K_1&=\operatorname{diag}
   ((\kappa/\kappa_*)(g_*/g),(g/g_*)(\eta/\eta_*)),\\
 K_2&=\operatorname{diag}(\eta/\eta_*,\kappa/\kappa_*).
 \end{aligned}                                               \tag{OP.3}
\]
Each determinant has winding \(-6\), and each \(K_i\) is the identity
at the initial parameter. Thus their real oriented normal loops have
even parity. We will not rely on parity alone: OP.2 gives their
explicit contractions.

Write the rho circle as \(\rho=\rho_*e^{i\theta}\). Each unit of
winding zero in (OP.3) has its actual periodic logarithm, chosen zero
at \(\theta=0\). For \(\kappa/\kappa_*\), retain the factor
\(e^{-6i\theta}\), and take the periodic logarithm of the remaining
unit. Multiplication by the exponentials of scalar multiples of these
logarithms supplies exact based homotopies of the unit factors. After
these carried unit maps, the remaining real normal loops are
\[
 K_0^{\rm ang}=K_2^{\rm ang}
        =\operatorname{diag}(1,e^{-6i\theta}),\qquad
 K_1^{\rm ang}=\operatorname{diag}(e^{-6i\theta},1).
                                                               \tag{OP.4}
\]
The same logarithms can be inserted in the final maps, as explained
in OP.3; the original units are not discarded.

\par

\medskip\noindent\textbf{OP.2. Explicit based quaternion contractions in the actual order}

\par\nobreak

For an ordered pair of complex normal coordinates \((n_0,n_1)\),
we use exactly the K23 convention
\[
                         Q=n_0+j\overline{n_1}.       \tag{OP.5}
\]
This real coordinate has the original complex normal orientation.
Choose a smooth function \(\vartheta:[0,1]\to[0,\pi/2]\), constant
near both ends and strictly increasing between them, and put
\[
 \begin{aligned}
 p(s,\theta)&=\sin\vartheta(s)\,e^{-3i\theta}
                                      +j\cos\vartheta(s),\\
 H(s,\theta)&=p(s,\theta)p(s,0)^{-1}.
 \end{aligned}                                               \tag{OP.6}
\]
Both factors are unit quaternions. The maps are smooth and periodic,
and satisfy
\[
 H(0,\theta)=1,\qquad H(1,\theta)=e^{-3i\theta},\qquad
 H(s,0)=H(s,2\pi)=1.                                  \tag{OP.7}
\]
Define two actual oriented real normal maps
\[
    T_-(s,\theta)(Q)=H QH^{-1},\qquad
    T_+(s,\theta)(Q)=H QH.                            \tag{OP.8}
\]
Their inverses are displayed by replacing the two factors with their
inverses in reverse order. They are orthogonal and orientation
preserving. At \(s=1\), the first sends the quaternion coefficients
\((n_0,\overline{n_1})\) to
\((n_0,e^{6i\theta}\overline{n_1})\). Thus in the original complex
coordinates it is
\((n_0,n_1)\mapsto(n_0,e^{-6i\theta}n_1)\).
The second is \((n_0,n_1)\mapsto(e^{-6i\theta}n_0,n_1)\).
These are precisely the two loops in (OP.4). In particular no
unconjugated quaternion coordinate is substituted for (OP.5).

We use \(T_-\) for \(C_0,C_2\) and \(T_+\) for \(C_1\).
Multiplying these normal maps by the diagonal exponentials of the
actual periodic unit logarithms, with a cutoff in \(s\), gives a
contraction \(\mathcal K_i(s,\theta)\) of each full loop (OP.3).
The cutoff is zero near \(s=0\) and one near \(s=1\), so
\[
 \mathcal K_i(0,\theta)=I,\quad
 \mathcal K_i(1,\theta)=K_i(\theta),\quad
 \mathcal K_i(s,0)=\mathcal K_i(s,2\pi)=I.             \tag{OP.9}
\]
Write these chosen contractions as
\(\mathcal K_i=\mathcal D_iT_i\), where \(\mathcal D_i(s,\theta)\)
is the actual diagonal exponential of the periodic unit logarithms,
and \(T_i=T_-\) for \(i=0,2\), \(T_i=T_+\) for \(i=1\).
All maps are invertible throughout. The diagonal unit homotopies
commute with the physical circle; the explicitly displayed maps
\(T_\pm\) need not do so.

\par

\medskip\noindent\textbf{OP.3. Gluing the actual ordinary plumbing product}

\par\nobreak

Form the marked regular neighbourhood from two small actual coordinate
balls at the vertices and three disjoint normal tubes along the curve
interiors. On each curve choose a radial coordinate \(s=s(|z|)\)
which is zero near its first vertex and one near its second. In the
first normal frame send
\[
                    (z,n)\longmapsto
                       (z,\mathcal K_i(s(|z|),\theta)n).
                                                               \tag{OP.10}
\]
It is the identity near the first vertex. Near the second it is
\(K_i\), and the equality
\[
               zD_i(\rho_*)K_i=zD_i(\rho)            \tag{OP.11}
\]
shows that it is exactly the second vertex comparison (OP.1).
Consequently these three maps and the two actual vertex maps agree
on their common collars. There is no unresolved corner: the open
curve tubes are disjoint away from the vertex balls, and (OP.11)
holds in the ordered normal directions of each incident axis.

The maps on the actual coordinate germs have the same first normal
derivatives as these displayed bundle maps. The remaining higher
terms are removed on smaller tube collars by normal rescaling.
For a map \(F\) preserving the zero section, put
\(R_t(z,n)=(z,tn)\) and use the literal conjugation
\(R_t^{-1}FR_t\) for \(t>0\). Its limit consists of the given
zero-section map and its invertible normal derivative. Only the
normal vector is scaled in this construction. On the compact parameter circle
one common sufficiently small tube makes this a smooth isotopy of
invertible germs through \(t=0\). On each vertex collar the original
unit maps have nonzero diagonal derivative and are treated by the
same rescaling, with the vertex coordinate balls protected.

For precision about gluing this isotopy, take its velocity
\((\partial_tF_t)F_t^{-1}\), multiply it by a smooth collar cutoff,
and integrate. This supplies the actual collar diffeomorphism
between the original germ and the displayed normal map. The cutoffs
are chosen on the three disjoint tube interiors and agree with the
already fixed vertex maps on their protected collars. Thus they
glue without changing any other incident axis. All isotopies are
periodic and are the identity at the reference parameter.

Finally transport the rounded boundary and use the resulting collar
isotopy to fix one regular-neighbourhood radius. This constructs a
periodic smooth product of the actual labelled plumbing with the
reference plumbing. The vertex maps are complex-unit maps, the
curve maps preserve their complex orientations, and the normal maps
are in the positive real group. Hence the product preserves the
six-manifold orientation. It is an ordinary smooth product of the
labelled neighbourhood; no source-height-preserving or
delta-equivariant product is asserted.

\par

\medskip\noindent\textbf{OP.4. The physical circle and the exact degree-three axis}

\par\nobreak

We now retain the physical circle in this product. After the actual
extending unit maps of OP.1, on the fixed-curve normal quaternion
(OP.5) it is \(Q\mapsto Qe^{-i\phi}\). Conjugating by the actual
map \(T_-\) gives
\[
       Q\longmapsto Q\,H(s,\theta)e^{-i\phi}H(s,\theta)^{-1}.
                                                               \tag{OP.12}
\]
The same right-hand axis occurs after multiplication by the fixed
source Hopf representative used to form \(q_H\). Thus the circle
axis in the ordinary product is
\[
                         \nu(s,\theta)=H iH^{-1}.     \tag{OP.13}
\]
It equals \(i\) on both ends of the cylinder, but it is not the
constant axis on its interior. The zero-winding diagonal unit maps
are the previously displayed transported maps; contracting their
logarithms carries (OP.12) to the formula written here and leaves
the following integer degree unchanged.

Orient the cylinder by \(ds\wedge d\theta\), collapse each boundary
circle to its endpoint, and orient the target sphere of imaginary
unit quaternions by its outward normal. We compute the degree rather
than infer it from parity. Put \(A=\sin\vartheta\),
\(B=\cos\vartheta\), \(C=\cos3\theta\), \(S=\sin3\theta\).
Then (OP.6) is literally
\[
 H=(B^2+A^2C)-A^2S\,i+AB(1-C)\,j+AB S\,k.            \tag{OP.14}
\]
The equation \(HiH^{-1}=-i\) is equivalent to zero real and
\(i\)-coefficients of \(H\). Its only solutions are
\[
 \vartheta=\pi/4,\qquad
 \theta=(\pi+2\pi j)/3,
                   \quad j=0,1,2.                    \tag{OP.15}
\]
At each one \(H=j\), and direct differentiation gives
\[
                    \partial_s\nu=4\vartheta'(s)k,
                    \partial_\theta\nu=3j.           \tag{OP.16}
\]
The ordered basis \((k,j)\) is positive at the target point \(-i\),
since \((-i,k,j)\) has the ambient orientation \((i,j,k)\).
Every preimage in (OP.15) is therefore regular and positive. Hence
\[
                              \deg\nu=+3.             \tag{OP.17}
\]

For completeness the two nonfixed normal actions are also explicit.
In their ordered normal coordinates of OP.1, the initial actions are
\[
 R_1(\phi)Q=e^{-i\phi/2}Qe^{-i\phi/2},\qquad
 R_2(\phi)Q=e^{i\phi/2}Qe^{-i\phi/2}.
\]
Their base coordinates have weights \(+1,-1\), respectively.
The ordinary product carries their normal actions to
\[
 \begin{aligned}
 C_1:\quad Q&\longmapsto
 (He^{-i\phi/2}H^{-1})Q(H^{-1}e^{-i\phi/2}H),\\
 C_2:\quad Q&\longmapsto
 (He^{i\phi/2}H^{-1})Q(He^{-i\phi/2}H^{-1}).
 \end{aligned}                                               \tag{OP.18}
\]
The paired central signs make both formulas close at \(\phi=2\pi\).
For the full original units, the physical normal action on each tube is
exactly
\[
 \mathcal D_i(s,\theta)
      [T_i(s,\theta)R_i(\phi)T_i(s,\theta)^{-1}]
                          \mathcal D_i(s,\theta)^{-1},
                                                               \tag{OP.19}
\]
with \(R_0(\phi)Q=Qe^{-i\phi}\). Thus the written \(H\)-axis and
(OP.18) are reached from the full-unit action by the displayed based
logarithm contraction, with the real normal metric transported along
that contraction. Together with their base rotations, (OP.19) gives
the actual transported circle actions on all three tubes.
They agree with the vertex action on the protected collars.

The product in OP.3 thus provides a possible ordinary common boundary
marking, with the nonconstant circle action (OP.12)--(OP.18). The
circle-trace and stable-derivative arguments for a constant equivariant
product require a further correction in this marking. That correction
is not evaluated here. In particular neither a rank-one initial
variation matrix nor vanishing nonfixed relative derivatives follows
from (OP.17).

\par

\medskip\noindent\textbf{OP.5. The closed fixed-curve frame comparison}

\par\nobreak

We also compute the actual normal-frame comparison on the fixed
curve. Write \(z=\tan\chi\,e^{i\psi}\), with
\(0\leq\chi\leq\pi/2\), and put
\[
 U(z)=\frac{1+jz}{\sqrt{1+|z|^2}},\qquad q_H=U(z)Q.
\]
The ordinary product of OP.2 and the physical Cox product of P3K.42
are therefore related by the literal arrow
\[
 q_H\longmapsto L(z,\theta)q_H R(z,\theta),\qquad
 L=UHU^{-1},\quad R=H^{-1}.                         \tag{OP.20}
\]
Here the radial parameter in \(H\) is the previously chosen smooth
function of \(|z|\). Both maps close on \(\mathbb P^1\times S^1\).
At the first pole \(H=1\). Near the second pole \(H=e^{-3i\theta}\),
so \(L=\cos3\theta-(UiU^{-1})\sin3\theta\) extends in the smooth
Hopf axis; its limiting value is \(e^{3i\theta}\). The right factor
extends independently of the source angle. The actual zero-winding
diagonal unit maps in OP.19 are carried in this comparison and then
contracted by their displayed based logarithms. This is a homotopy
of the actual closed frame arrow, rather than a discarded unit.

We orient the domain by its complex curve orientation followed by
positive \(\theta\), hence by \(d\chi\wedge d\psi\wedge d\theta\)
on the cylinder. In the quaternion orientation \((1,i,j,k)\), use
the matrix trace with \(\operatorname{tr}(ijk)=-2\), and the positive
normalized form
\[
 \omega(g)=-\frac{\operatorname{tr}(g^{-1}dg)^3}{24\pi^2}.
\]
The Maurer--Cartan equations give the exact product identity
\[
 \omega(gh)=\omega(g)+\omega(h)+
 \frac{1}{8\pi^2}d\operatorname{tr}
       (g^{-1}dg\wedge dh\,h^{-1}).                  \tag{OP.21}
\]
Applying it twice to \(UHU^{-1}\), the two \(U\) cubic terms cancel
and the \(H\) term vanishes because \(H\) depends on only two
variables. Thus its integral is the boundary integral of
\[
 \frac{1}{8\pi^2}\left\{
 \operatorname{tr}(U^{-1}dU\wedge dH\,H^{-1})+
 \operatorname{tr}((UH)^{-1}d(UH)\wedge dU^{-1}U)\right\}.
                                                               \tag{OP.22}
\]
There is no first-pole contribution. At the second pole,
\(U=je^{i\psi}\), \(H=e^{-3i\theta}\), and
\[
 U^{-1}dU=i\,d\psi,\qquad dH\,H^{-1}=-3i\,d\theta.
\]
Each trace in braces in (OP.22) is
\(6\,d\psi\wedge d\theta\). The outward boundary orientation there
is \(d\psi\wedge d\theta\), and consequently
\[
 \deg L=\frac{12(2\pi)^2}{8\pi^2}=6,\qquad
 \deg R=0.                                            \tag{OP.23}
\]
The second equality follows directly from its two-variable
factorization. The Hopf coordinate \(h=UiU^{-1}\) reverses the
complex orientation of the source curve: near \(z=0\), its two
derivatives in \((\Re z,\Im z)\) are \((-2k,-2j)\), negatively
oriented at \(i\). Therefore the same closed comparison has
factor degrees \((-6,0)\) in the \(h\)-sphere-times-parameter
orientation of P3K.23, and \((6,0)\) in complex-curve orientation.
We have computed these closed factor maps and their degrees. The
connecting-homomorphism sign and the resulting relative-return
correction remain unevaluated; no additional fan return follows
without those arrows.

The exact normal-matrix, quaternion and degree checks are in
\path{code/part3/source_cusp_ordinary_product.py}.
\end{proof}